Generalized disjunctive programming (GDP) describes process design alternatives with Boolean variables, disjunctions, and logical propositions, and fixing one configuration leaves a subproblem in the continuous variables.
Conventional GDP methods require explicit algebraic formulations that simulators may not expose.
Even when such formulations are available, nonconvex models can be intractable.
The external variables define an integer lattice that derivative-free optimization (DFO) searches using subproblem values.
We compare surrogate DFO methods and the direct-search method \texttt{NOMAD} with discrete steepest descent (\texttt{D-SDA}), logic-based solvers, mixed-integer nonlinear programming (MINLP) solvers, and enumeration.
The benchmarks include nine process models, 21 test functions, a pharmaceutical flowsheet simulated in \texttt{PharmaPy}, and a catalytic distillation column with 100,947 configurations.
On the largest reactor network, the surrogate methods \texttt{MISO} and \texttt{RBFOpt} reach the reference in nine or ten of ten runs with about 30 subproblem solves, whereas logic-based methods solve nearly one subproblem for each admissible configuration and most \texttt{D-SDA} descents end at one local minimum.
In the flowsheet example, where GDP solvers do not apply, \texttt{RBFOpt} and \texttt{MISO} find the best of 36 simulated configurations in every run with 20 evaluations.
For the catalytic column, \texttt{D-SDA} and \texttt{MISO} reach the best known design with 100 evaluations, whereas MINLP solvers find no design within 15\% of it and global solvers return contradictory certificates.
This search pays off when the configuration space is too large to enumerate, has deceptive local minima, or has MINLP reformulations that fail, and enumeration or a direct reformulation is simpler for small algebraic models.
